\documentclass[12pt]{report}
\input{preamble.tex}
\usepackage{float}
\usepackage{csquotes}
\usepackage{translator}
\uselanguage{Spanish}
\languagepath{Spanish}
\usepackage{siunitx}
\usepackage{pgfgantt}

%\usepackage{draftwatermark}
%\SetWatermarkLightness{0.9}

\makeatletter
\newcommand*{\centerfloat}{%
  \parindent \z@
  \leftskip \z@ \@plus 1fil \@minus \textwidth
  \rightskip\leftskip
  \parfillskip \z@skip}
\makeatother

\graphicspath{{./figs/}}
% Renombro lo capitulos
%\addto\captionsspanish{\renewcommand{\chaptername}{Clase}}

\hypersetup{pdftitle={TITULO METADATO - versión XXXX},pdfauthor={AUTORES METADATO}}


\title{\vfill{\large SUBTITULO} \\ TITULO \\ {\large TIPO DE DOCUMENTO - versión XXX}}
\author[*]{AUTOR 1}

\author[*]{AUTOR 2}

\author[*]{AUTOR 3}

\affil[*]{Unidad de Formación Masiva}

\titlepic{\vfill \includegraphics[height=1.5cm]{../figs/logo-conae.png}\hfill\includegraphics[height=1.5cm]{../figs/logo-ministerio.png}}

\date{}
\setcounter{tocdepth}{1}
\begin{document}
%\frontmatter
\maketitle
\titlepage
%\tableofcontents
%\include{prefacio}
\addtolength{\voffset}{-2cm}
\addtolength{\headsep}{0.5cm}
\addtolength{\textheight}{1cm}
\newpage
\tableofcontents
\input{revisor}

%\mainmatter
\chapter{General}

LINEAMIENTOS GENERALES

El curso tiene como objetivo:

\begin{displayquote}
\emph{OBJETIVO CON EMFASIS}
\end{displayquote}

DESTINATARIOS

\begin{itemize}
    \item DESTINATARIO 1
    \item DESTINATARIO 2
    \item DESTINATARIO 3
\end{itemize}

DESCRIPCION CURSO

\section{Requisitos del curso}

Para realizar el curso los participantes deben:

\begin{itemize}
    \item Haber realizado y aprobado el \href{https://www.argentina.gob.ar/cursos-y-capacitaciones/nivel-i-introduccion-la-teledeteccion}{Curso Nivel I: Introducción a Teledetección} o tener conocimientos equivalentes. 
    \item Contar con una PC con Windows 7 o superior, Linux. Procesador i3 o superior, 4GB de memoria RAM, 30GB de espacio en disco disponible.
    \item Desenvolverse con soltura en un entorno de trabajo informático: crear una carpeta, localizar un archivo en la PC, exportar archivos a formato de imágenes y pdf, descomprimir un archivo e instalar archivos ejecutables.
    \item Disponer de una conexión a internet de 1 Mbit/s.
\end{itemize}


\section{Programa}

El curso incluye los siguientes temas:

\begin{itemize}
    \item CLASE 1

    \item CLASE 2
    
    \item CLASE 3

    \item CLASE 4

    \item CLASE 5

\end{itemize}

\section{Software}

SNAP 7. PORQUE.


\section{Cronograma}
El cronograma de diseño e implementación es el siguiente:

\begin{ganttchart}[x unit=0.04cm, y unit chart = .7cm,
                    y unit title = .7cm,
    time slot format=isodate,
    calendar week text = {\currentweek{}},canvas/.append style={fill=none},
    link bulge=5, link tolerance = 10]{2019-12-15}{2020-11-30}
\gantttitlecalendar{year, month=shortname}\\
\ganttgroup{Diseño}{2019-12-16}{2020-01-24}\\
\ganttgroup{Armado}{2020-01-24}{2020-07-31}\\
\ganttbar{Práctica}{2020-01-24}{2020-04-17}\\
\ganttbar{Teoría}{2020-04-03}{2020-06-12}\\
\ganttbar{Plataforma}{2020-04-03}{2020-07-31}\\
\ganttbar{Evaluación}{2020-05-15}{2020-06-26}\\
\ganttgroup{Dictado}{2020-9-7}{2020-11-30}
\begin{scope}[on background layer]
\ganttlink{elem0}{elem1}
\ganttlink{elem1}{elem6}
\end{scope}
\end{ganttchart}

\subsection{Listado de tareas}

% \usepackage{booktabs}
\begin{table}[h!]
\centerfloat
\small
\begin{tabular}{@{}llclc@{}}
\toprule
Grupo   & Tarea                & Código & Responsable             & Entrega      \\ \midrule
Diseño  & Documento de diseño  & \href{https://www.overleaf.com/}{1.1}    & Todos                   & 2020-01-24 \\ \midrule
Armado  & Práctica 1           & \href{https://www.overleaf.com/}{2.1}    & QUIEN      & 2020-02-14 \\
        & Práctica 2           & \href{https://www.overleaf.com/}{2.2}    & QUIEN        & 2020-02-14 \\
        & Práctica 3           & \href{https://www.overleaf.com/}{2.3}    & QUIEN        & 2020-02-28 \\
        & Práctica 4           & \href{https://www.overleaf.com/}{2.4}    & QUIEN      & 2020-03-13 \\
        & Práctica 5           & \href{https://www.overleaf.com/}{2.5}    & QUIEN      & 2020-04-17 \\
        & Teórica 1            & \href{https://www.overleaf.com/}{2.6}    & QUIEN        & 2020-04-17 \\
        & Teórica 2            & \href{https://www.overleaf.com/}{2.7}    & QUIEN        & 2020-05-01 \\
        & Teórica 3            & \href{https://www.overleaf.com/}{2.8}    & QUIEN        & 2020-05-15 \\
        & Teórica 4            & \href{https://www.overleaf.com/}{2.9}    & QUIEN      & 2020-05-01 \\
        & Teórica 5            & \href{https://www.overleaf.com/}{2.10}   & QUIEN      & 2020-05-15 \\
        & Audio videos         & 2.11   & Todos                   & 2020-05-29 \\
        & Edición videos       & 2.12   & QUIEN       & 2020-06-12 \\
        & Armado plataforma    & 2.13   & QUIEN        & 2020-06-05 \\
        & Evaluación 1         & 2.14   & QUIEN      & 2020-06-05 \\
        & Evaluación 2         & 2.15   & QUIEN      & 2020-06-12 \\
        & Evaluación 3         & 2.16   & QUIEN      & 2020-06-19 \\
        & Evaluación 4         & 2.17   & QUIEN      & 2020-06-26 \\
        & Revisión general     & 2.18   & QUIEN             & 2020-07-31 \\\midrule
Dictado & Dictado del curso    & 3.1    & Todos                   & 2020-09-07 \\ 
        & Evaluación del curso & 3.2    & QUIEN        & 2020-11-30 \\
        & Revisión del curso   & 3.3    & Todos                   & 2020-11-30 \\ \bottomrule
\end{tabular}
\caption{Listado de tareas del curso.}
\end{table}

\chapter{Teoría}\label{teoria}

Se detallan cuales son las temáticas para cada clase que se incluyen en los videos teóricos y, en casos específicos, se incluye un documento con información adicional.

\section{Clase 1}
Video 

\begin{itemize}
    \item TEMA 1
    \item TEMA 2
    \item TEMA 3
    \item TEMA 4
    \item TEMA 5
\end{itemize}

Documento adicional: COSAS COPADAS

\section{Clase 2}

Video

\begin{itemize}
    \item TEMA 1
    \item TEMA 2
    \item TEMA 3
    \item TEMA 4
    \item TEMA 5
\end{itemize}

Documento adicional: COSAS COPADAS

\section{Clase 3}

Video

\begin{itemize}
    \item TEMA 1
    \item TEMA 2
    \item TEMA 3
    \item TEMA 4
    \item TEMA 5
\end{itemize}

Documento adicional: COSAS COPADAS

\section{Clase 4}

\begin{itemize}
    \item TEMA 1
    \item TEMA 2
    \item TEMA 3
    \item TEMA 4
    \item TEMA 5
\end{itemize}

Documento adicional: COSAS COPADAS

\section{Clase 5}

\begin{itemize}
    \item TEMA 1
    \item TEMA 2
    \item TEMA 3
    \item TEMA 4
    \item TEMA 5
\end{itemize}

Documento adicional: COSAS COPADAS

\chapter{Práctica}\label{practica}

A continuación se describe las imágenes y herramientas utilizadas en cada guía práctica.

El material del curso se encuentra en

\begin{center}
    \href{https://cloud.conae.gov.ar}{https://cloud.conae.gov.ar}
\end{center}

con contraseña

\begin{center}
    \texttt{password01}
\end{center}

\section{Clase 1}
DESCRIPCION

Las herramientas usadas son:

\begin{itemize}
    \item \menu{Herramienta 1} - Herramienta 1
    \item \menu{Herramienta 2} - Herramienta 2
\end{itemize}

El esquema general de la clase es:

\begin{enumerate}
    \item Paso 1
    \item Paso 2
    \item Paso 3
\end{enumerate}

Los productos son 

\begin{itemize}
    \item \directory{IMAGEN1.tif}
\end{itemize}

\section{Clase 2}

DESCRIPCION

Las herramientas usadas son:

\begin{itemize}
    \item \menu{Herramienta 1} - Herramienta 1
    \item \menu{Herramienta 2} - Herramienta 2
\end{itemize}

El esquema general de la clase es:

\begin{enumerate}
    \item Paso 1
    \item Paso 2
    \item Paso 3
\end{enumerate}

Los productos son 

\begin{itemize}
    \item \directory{IMAGEN1.tif}
\end{itemize}

\section{Clase 3}
DESCRIPCION

Las herramientas usadas son:

\begin{itemize}
    \item \menu{Herramienta 1} - Herramienta 1
    \item \menu{Herramienta 2} - Herramienta 2
\end{itemize}

El esquema general de la clase es:

\begin{enumerate}
    \item Paso 1
    \item Paso 2
    \item Paso 3
\end{enumerate}

Los productos son 

\begin{itemize}
    \item \directory{IMAGEN1.tif}
\end{itemize}

\section{Clase 4}
DESCRIPCION

Las herramientas usadas son:

\begin{itemize}
    \item \menu{Herramienta 1} - Herramienta 1
    \item \menu{Herramienta 2} - Herramienta 2
\end{itemize}

El esquema general de la clase es:

\begin{enumerate}
    \item Paso 1
    \item Paso 2
    \item Paso 3
\end{enumerate}

Los productos son 

\begin{itemize}
    \item \directory{IMAGEN1.tif}
\end{itemize}

\section{Clase 5}

DESCRIPCION

Las herramientas usadas son:

\begin{itemize}
    \item \menu{Herramienta 1} - Herramienta 1
    \item \menu{Herramienta 2} - Herramienta 2
\end{itemize}

El esquema general de la clase es:

\begin{enumerate}
    \item Paso 1
    \item Paso 2
    \item Paso 3
\end{enumerate}

Los productos son 

\begin{itemize}
    \item \directory{IMAGEN1.tif}
\end{itemize}


\chapter{Evaluación}\label{evaluacion}

La evaluación del curso se realiza a través de preguntas de autocorrección para las clases 1, 2, 3 y 4, mientras que en la clase 5 se pide un mapa en formato .jpg, que corrige el tutor a cargo. Todas las actividades de evaluación se desarrollan a través de la Plataforma de Educación a Distancia de la Unidad de Formación Masiva.

\section{Clase 1}

En la clase 1 se evalúa:

\begin{itemize}
    \item PREGUNTA 1
    \item PREGUNTA 2
\end{itemize}

\section{Clase 2}

En la clase 2 se evalúa

\begin{itemize}
    \item PREGUNTA 1
    \item PREGUNTA 2
\end{itemize}

\section{Clase 3}

En la clase 3 se evalúa

\begin{itemize}
    \item PREGUNTA 1
    \item PREGUNTA 2
\end{itemize}

\section{Clase 4}

En la clase 3 se evalúa:

\begin{itemize}
    \item PREGUNTA 1
    \item PREGUNTA 2
\end{itemize}

\section{Clase 5}

TAREA 1

\chapter{Actividades adicionales}\label{adicional}

\section{Clase 1}

ADICIONAL 1

\section{Clase 2}

ADICIONAL 1

\section{Clase 3}

ADICIONAL 1
\section{Clase 4}

ADICIONAL 1


\section{Clase 5}

No corresponde.

\nocite{*}
 
\printbibliography

\appendix
%\include{apendice1}
%\include{apendice2}

%\backmatter

\end{document}
